\documentclass[11pt]{article}

\usepackage[a4paper,margin=1in]{geometry}
\usepackage{amsmath,amssymb,mathtools}
\usepackage{booktabs}
\usepackage{hyperref}

\newcommand{\dd}{\mathrm d}
\newcommand{\ii}{\mathrm i}
\newcommand{\eps}{\epsilon}
\newcommand{\U}{\mathcal U}
\newcommand{\F}{\mathcal F}
\newcommand{\V}{\mathcal V}
\newcommand{\hU}{\widehat{\mathcal U}}
\newcommand{\hF}{\widehat{\mathcal F}}
\newcommand{\Dlt}{\Delta}

\title{Fast Sector Decomposition from Black-Box Symanzik Polynomials}
\author{}
\date{}

\begin{document}

\maketitle

\section{The example triangle}

We consider the scalar one-loop triangle with one incoming off-shell momentum
and two outgoing massless momenta,
\begin{equation}
  p_0=p_1+p_2,
  \qquad
  p_1^2=p_2^2=0,
  \qquad
  p_0^2=s=m_0^2 .
\end{equation}
All three internal propagators are taken to have the same mass \(m\).  A
convenient routing is
\begin{align}
  D_0 &= \ell^2-m^2+\ii0,\\
  D_1 &= (\ell+p_1)^2-m^2+\ii0,\\
  D_2 &= (\ell-p_2)^2-m^2+\ii0 .
\end{align}
The scalar integral is
\begin{equation}
  C_0(s;m^2)
  =
  \int \frac{\dd^D\ell}{\ii\pi^{D/2}}\,
  \frac{1}{D_0D_1D_2}.
\end{equation}
With Feynman parameters \(x_i\geq 0\), the projective representation is
\begin{equation}
  C_0(s;m^2)
  =
  -\Gamma(\alpha)
  \int_0^\infty \dd x_0\,\dd x_1\,\dd x_2\,
  \delta(1-x_0-x_1-x_2)\,
  \frac{\U^{3-D}}{\F^\alpha},
  \qquad
  \alpha=3-\frac{D}{2}.
\end{equation}
For this routing,
\begin{equation}
  \U=x_0+x_1+x_2,
\end{equation}
and
\begin{equation}
  \F
  =
  m^2\U^2
  -
  s\,x_1x_2
  -
  \ii0 .
\end{equation}
Solving the simplex constraint by
\begin{equation}
  x_0=1-x_1-x_2
\end{equation}
gives
\begin{equation}
  0\leq x_1\leq 1,
  \qquad
  0\leq x_2\leq 1-x_1,
\end{equation}
and therefore
\begin{equation}
  \U=1,
  \qquad
  \F=m^2-s\,x_1x_2-\ii0 .
\end{equation}
For \(D=4-2\eps\), one has \(\alpha=1+\eps\), and
\begin{equation}
  C_0(s;m^2)
  =
  -\Gamma(1+\eps)
  \int_0^1 \dd x_1
  \int_0^{1-x_1} \dd x_2\,
  \left(m^2-s\,x_1x_2-\ii0\right)^{-1-\eps}.
\end{equation}
Throughout the rest of the note, sector integrands \(I_s\) do not include the
global prefactor \(-\Gamma(\alpha)\).  Thus an exact decomposition has the form
\begin{equation}
  C_0(s;m^2)
  =
  -\Gamma(\alpha)
  \sum_s
  \int_{\mathcal C_s}\dd y_s\,
  I_s(y_s).
\end{equation}

\subsection{Endpoint sectors}

The primary endpoint decomposition splits the simplex into two sectors.

\paragraph{Sector \(S_1\), \(x_1\geq x_2\).}
\begin{equation}
  X_{S_1}(t,z)
  =
  (x_0,x_1,x_2)
  =
  \left(
    1-t,\,
    \frac{t}{1+z},\,
    \frac{tz}{1+z}
  \right),
  \qquad
  0\leq t,z\leq 1 .
\end{equation}
The Jacobian is
\begin{equation}
  J_{S_1}(t,z)=\frac{t}{(1+z)^2}.
\end{equation}
The pulled-back second Symanzik polynomial is
\begin{equation}
  \F_{S_1}(t,z)
  =
  m^2
  -
  s\,\frac{t^2z}{(1+z)^2}
  -
  \ii0 .
\end{equation}
In the massless endpoint case \(m=0\),
\begin{equation}
  \F_{S_1}(t,z)
  =
  (-s-\ii0)\,\frac{t^2z}{(1+z)^2}.
\end{equation}
The expected monomial is therefore
\begin{equation}
  M_{S_1}(t,z)=t^2z .
\end{equation}

\paragraph{Sector \(S_2\), \(x_2\geq x_1\).}
\begin{equation}
  X_{S_2}(t,z)
  =
  (x_0,x_1,x_2)
  =
  \left(
    1-t,\,
    \frac{tz}{1+z},\,
    \frac{t}{1+z}
  \right),
  \qquad
  0\leq t,z\leq 1 .
\end{equation}
The Jacobian is again
\begin{equation}
  J_{S_2}(t,z)=\frac{t}{(1+z)^2}.
\end{equation}
The pulled-back polynomial is
\begin{equation}
  \F_{S_2}(t,z)
  =
  m^2
  -
  s\,\frac{t^2z}{(1+z)^2}
  -
  \ii0 .
\end{equation}
In the massless endpoint case,
\begin{equation}
  \F_{S_2}(t,z)
  =
  (-s-\ii0)\,\frac{t^2z}{(1+z)^2},
  \qquad
  M_{S_2}(t,z)=t^2z .
\end{equation}
With these definitions, the two-sector representation of the original
triangle is
\begin{align}
  C_0(s;m^2)
  &=
  -\Gamma(\alpha)
  \sum_{a=1}^{2}
  \int_0^1\dd t
  \int_0^1\dd z\,
  I_{S_a}(t,z),
  \\
  I_{S_a}(t,z)
  &=
  \frac{t}{(1+z)^2}\,
  \U(X_{S_a}(t,z))^{3-D}\,
  \F(X_{S_a}(t,z))^{-\alpha}.
\end{align}
Equivalently, after inserting the explicit polynomial,
\begin{equation}
  C_0(s;m^2)
  =
  -\Gamma(\alpha)
  \sum_{a=1}^{2}
  \int_0^1\dd t
  \int_0^1\dd z\,
  \frac{t}{(1+z)^2}
  \left(
    m^2
    -
    s\,\frac{t^2z}{(1+z)^2}
    -
    \ii0
  \right)^{-\alpha}.
\end{equation}
For \(D=4-2\eps\) and \(m=0\), this becomes
\begin{equation}
  C_0(s;0)
  =
  -\Gamma(1+\eps)
  \sum_{a=1}^{2}
  \int_0^1\dd t
  \int_0^1\dd z\,
  t^{-1-2\eps}
  z^{-1-\eps}
  G_{S_a}(t,z),
\end{equation}
where
\begin{equation}
  G_{S_a}(t,z)
  =
  (-s-\ii0)^{-1-\eps}(1+z)^{2\eps}.
\end{equation}

\section{Direct approach}

In the direct approach, each sector map \(X_s\) is substituted explicitly into
\(\U\) and \(\F\).  If the pulled-back polynomial factorizes as
\begin{equation}
  \F_s(y)
  =
  \F(X_s(y))
  =
  M_s(y)\,\phi_s(y),
\end{equation}
then the sector integrand is built as
\begin{equation}
  I_s^{\mathrm{dir}}(y)
  =
  J_s(y)\,
  \U(X_s(y))^{3-D}\,
  M_s(y)^{-\alpha}\,
  \phi_s(y)^{-\alpha}.
\end{equation}
The corresponding contribution to the original scalar integral is always
understood as
\begin{equation}
  C_0
  =
  -\Gamma(\alpha)
  \sum_s
  \int_{\mathcal C_s}\dd y_s\,
  I_s^{\mathrm{dir}}(y_s),
\end{equation}
with the sector domains \(\mathcal C_s\) specified explicitly below.

\subsection{\texorpdfstring{Endpoint sector \(S_1\)}{Endpoint sector S1}}

For \(m=0\),
\begin{equation}
  \U_{S_1}(t,z)=1,
\end{equation}
and
\begin{equation}
  \F_{S_1}(t,z)
  =
  (-s-\ii0)\,\frac{t^2z}{(1+z)^2}
  =
  M_{S_1}(t,z)\,\phi_{S_1}(t,z),
\end{equation}
with
\begin{equation}
  M_{S_1}(t,z)=t^2z,
  \qquad
  \phi_{S_1}(t,z)=\frac{-s-\ii0}{(1+z)^2}.
\end{equation}
Therefore
\begin{align}
  I_{S_1}^{\mathrm{dir}}(t,z)
  &=
  \frac{t}{(1+z)^2}\,
  (t^2z)^{-\alpha}
  \left(\frac{-s-\ii0}{(1+z)^2}\right)^{-\alpha}.
\end{align}
For \(D=4-2\eps\), i.e. \(\alpha=1+\eps\),
\begin{equation}
  I_{S_1}^{\mathrm{dir}}(t,z)
  =
  (-s-\ii0)^{-1-\eps}\,
  t^{-1-2\eps}\,
  z^{-1-\eps}\,
  (1+z)^{2\eps}.
\end{equation}

\subsection{\texorpdfstring{Endpoint sector \(S_2\)}{Endpoint sector S2}}

By symmetry,
\begin{equation}
  \F_{S_2}(t,z)
  =
  (-s-\ii0)\,\frac{t^2z}{(1+z)^2},
  \qquad
  M_{S_2}(t,z)=t^2z,
\end{equation}
and hence
\begin{equation}
  I_{S_2}^{\mathrm{dir}}(t,z)
  =
  (-s-\ii0)^{-1-\eps}\,
  t^{-1-2\eps}\,
  z^{-1-\eps}\,
  (1+z)^{2\eps}.
\end{equation}
The full massless triangle is therefore
\begin{align}
  C_0(s;0)
  &=
  -\Gamma(1+\eps)
  \sum_{a=1}^{2}
  \int_0^1\dd t
  \int_0^1\dd z\,
  I_{S_a}^{\mathrm{dir}}(t,z)
  \notag\\
  &=
  -2\Gamma(1+\eps)
  (-s-\ii0)^{-1-\eps}
  \int_0^1\dd t
  \int_0^1\dd z\,
  t^{-1-2\eps}
  z^{-1-\eps}
  (1+z)^{2\eps}.
\end{align}

\section{Indirect approach}

The indirect approach keeps the Symanzik polynomials closed.  We assume only
black-box evaluators
\begin{equation}
  \hU(x_0,x_1,x_2),
  \qquad
  \hF(x_0,x_1,x_2),
\end{equation}
together with their partial derivatives with respect to the original
parameters \(x_i\).  In practice, these derivative jets can be computed
efficiently with dual numbers.

For a sector map \(X_s\), define the pullbacks
\begin{equation}
  U_s(y)=\hU(X_s(y)),
  \qquad
  F_s(y)=\hF(X_s(y)).
\end{equation}
Derivatives with respect to sector variables are obtained by chain rules.  For
example,
\begin{equation}
  \partial_a F_s(y)
  =
  \sum_i
  \frac{\partial \hF}{\partial x_i}(X_s(y))\,
  \frac{\partial X_{s,i}}{\partial y_a}(y),
\end{equation}
and
\begin{align}
  \partial_a\partial_b F_s(y)
  &=
  \sum_{i,j}
  \frac{\partial^2 \hF}{\partial x_i\partial x_j}(X_s(y))\,
  \frac{\partial X_{s,i}}{\partial y_a}(y)\,
  \frac{\partial X_{s,j}}{\partial y_b}(y)
  \notag\\
  &\quad
  +
  \sum_i
  \frac{\partial \hF}{\partial x_i}(X_s(y))\,
  \frac{\partial^2 X_{s,i}}{\partial y_a\partial y_b}(y).
\end{align}
Higher derivatives are obtained similarly, or more directly by propagating
dual-number jets through the known sector map and then through the black-box
polynomial evaluator.

The object to construct is the regular coefficient \(\phi_s\) defined by
\begin{equation}
  F_s(y)=M_s(y)\,\phi_s(y),
\end{equation}
where \(M_s\) is the expected monomial for that sector.  Once \(\phi_s\) is
available, the sector integrand is
\begin{equation}
  I_s^{\mathrm{ind}}(y)
  =
  J_s(y)\,
  U_s(y)^{3-D}\,
  M_s(y)^{-\alpha}\,
  \phi_s(y)^{-\alpha}.
\end{equation}
At interior points this is pointwise identical to direct substitution, since
\begin{equation}
  \phi_s(y)=\frac{\hF(X_s(y))}{M_s(y)}
\end{equation}
there.
For the endpoint decomposition before any monomial extraction, the black-box
form of the exact triangle is
\begin{align}
  C_0(s;m^2)
  &=
  -\Gamma(\alpha)
  \sum_{a=1}^{2}
  \int_0^1\dd t
  \int_0^1\dd z\,
  \frac{t}{(1+z)^2}
  \hU(X_{S_a}(t,z))^{3-D}
  \hF(X_{S_a}(t,z))^{-\alpha}.
\end{align}
When \(m=0\), and hence
\(\hF(X_{S_a}(t,z))=t^2z\,\phi_{S_a}(t,z)\), this same equality becomes
\begin{align}
  C_0(s;0)
  &=
  -\Gamma(1+\eps)
  \sum_{a=1}^{2}
  \int_0^1\dd t
  \int_0^1\dd z\,
  \frac{t}{(1+z)^2}\,
  \hU(X_{S_a}(t,z))^{-1+2\eps}
  \notag\\
  &\hspace{3.2cm}\times
  (t^2z)^{-1-\eps}\,
  \phi_{S_a}(t,z)^{-1-\eps}.
\end{align}

\subsection{Using subtraction terms}

The subtraction construction is the one used to expose the Laurent poles in
\(\eps\).  At interior points one evaluates
\begin{equation}
  \phi_s(y)=\frac{\hF(X_s(y))}{M_s(y)}.
\end{equation}
Boundary values are supplied by the continuous extension, using derivatives.
For the first endpoint sector,
\begin{align}
  \phi_{S_1}(0,z)
  &=
  \frac{1}{2z}
  \left[
    \partial_T^2 F_{S_1}(T,z)
  \right]_{T=0},
  \\
  \phi_{S_1}(t,0)
  &=
  \frac{1}{t^2}
  \left[
    \partial_Z F_{S_1}(t,Z)
  \right]_{Z=0},
  \\
  \phi_{S_1}(0,0)
  &=
  \frac12
  \left[
    \partial_T^2\partial_Z F_{S_1}(T,Z)
  \right]_{T=0,\;Z=0}.
\end{align}
Again, the derivatives of \(F_{S_1}\) are not symbolic derivatives of an opened
polynomial.  They are derivatives of \(\hF\circ X_{S_1}\), obtained from the
black-box derivative jets and the chain rule.

For \(D=4-2\eps\), write each factorized endpoint integrand as
\begin{equation}
  I_{S_a}^{\mathrm{ind}}(t,z)
  =
  t^{-1-2\eps}\,
  z^{-1-\eps}\,
  G_{S_a}(t,z),
  \qquad
  a=1,2,
\end{equation}
where
\begin{equation}
  G_{S_a}(t,z)
  =
  \frac{1}{(1+z)^2}\,
  U_{S_a}(t,z)^{-1+2\eps}\,
  \phi_{S_a}(t,z)^{-1-\eps}.
\end{equation}
The two-dimensional endpoint subtraction identity is
\begin{align}
  &\int_0^1 \dd t\int_0^1 \dd z\,
  t^{-1-2\eps}z^{-1-\eps}G(t,z)
  \notag\\
  &=
  \int_0^1 \dd t\int_0^1 \dd z\,
  t^{-1-2\eps}z^{-1-\eps}
  \left[
    G(t,z)-G(0,z)-G(t,0)+G(0,0)
  \right]
  \notag\\
  &\quad
  -\frac{1}{2\eps}
  \int_0^1 \dd z\,
  z^{-1-\eps}
  \left[
    G(0,z)-G(0,0)
  \right]
  \notag\\
  &\quad
  -\frac{1}{\eps}
  \int_0^1 \dd t\,
  t^{-1-2\eps}
  \left[
    G(t,0)-G(0,0)
  \right]
  +
  \frac{G(0,0)}{2\eps^2}.
\end{align}
Each bracket is less singular at the corresponding endpoint.  This is the
mechanism that turns endpoint singularities into explicit poles and finite
integrals.
Applied sector by sector, it gives the full massless triangle as
\begin{align}
  C_0(s;0)
  &=
  -\Gamma(1+\eps)
  \sum_{a=1}^{2}
  \Bigg\{
  \int_0^1 \dd t\int_0^1 \dd z\,
  t^{-1-2\eps}z^{-1-\eps}
  \Big[
    G_{S_a}(t,z)-G_{S_a}(0,z)
  \notag\\
  &\hspace{5.1cm}
    -G_{S_a}(t,0)+G_{S_a}(0,0)
  \Big]
  \notag\\
  &\qquad
  -\frac{1}{2\eps}
  \int_0^1 \dd z\,
  z^{-1-\eps}
  \left[
    G_{S_a}(0,z)-G_{S_a}(0,0)
  \right]
  \notag\\
  &\qquad
  -\frac{1}{\eps}
  \int_0^1 \dd t\,
  t^{-1-2\eps}
  \left[
    G_{S_a}(t,0)-G_{S_a}(0,0)
  \right]
  +
  \frac{G_{S_a}(0,0)}{2\eps^2}
  \Bigg\}.
\end{align}

For the explicit massless triangle in sector \(S_1\),
\begin{equation}
  U_{S_1}(t,z)=1,
  \qquad
  \phi_{S_1}(t,z)=\frac{-s-\ii0}{(1+z)^2}.
\end{equation}
Thus
\begin{equation}
  G_{S_1}^{\mathrm{exp}}(t,z)
  =
  (-s-\ii0)^{-1-\eps}(1+z)^{2\eps}.
\end{equation}
This function is independent of \(t\).  Therefore
\begin{equation}
  G_{S_1}^{\mathrm{exp}}(t,z)
  -
  G_{S_1}^{\mathrm{exp}}(0,z)
  -
  G_{S_1}^{\mathrm{exp}}(t,0)
  +
  G_{S_1}^{\mathrm{exp}}(0,0)
  =
  0,
\end{equation}
and
\begin{equation}
  G_{S_1}^{\mathrm{exp}}(t,0)
  -
  G_{S_1}^{\mathrm{exp}}(0,0)
  =
  0.
\end{equation}
The only non-trivial subtraction piece is
\begin{equation}
  G_{S_1}^{\mathrm{exp}}(0,z)
  -
  G_{S_1}^{\mathrm{exp}}(0,0)
  =
  (-s-\ii0)^{-1-\eps}
  \left[
    (1+z)^{2\eps}-1
  \right].
\end{equation}
The explicitly subtracted result for one endpoint sector is therefore
\begin{align}
  &\int_0^1 \dd t\int_0^1 \dd z\,
  I_{S_1}^{\mathrm{dir}}(t,z)
  \notag\\
  &=
  -\frac{(-s-\ii0)^{-1-\eps}}{2\eps}
  \int_0^1 \dd z\,
  z^{-1-\eps}
  \left[
    (1+z)^{2\eps}-1
  \right]
  +
  \frac{(-s-\ii0)^{-1-\eps}}{2\eps^2}.
\end{align}
The second sector gives the same contribution, so the full massless triangle is
\begin{align}
  C_0(s;0)
  &=
  -2\Gamma(1+\eps)
  \left\{
  -\frac{(-s-\ii0)^{-1-\eps}}{2\eps}
  \int_0^1 \dd z\,
  z^{-1-\eps}
  \left[
    (1+z)^{2\eps}-1
  \right]
  +
  \frac{(-s-\ii0)^{-1-\eps}}{2\eps^2}
  \right\}.
\end{align}
This expression was obtained without requiring any symbolic expansion in an
implementation: only \(X_{S_a}\), \(J_{S_a}\), the known monomial \(t^2z\), and
black-box values and derivative limits of \(\hU\) and \(\hF\) are required in
each endpoint sector.

\section{General parametric structure}

For a scalar \(L\)-loop integral with \(N\) propagators of powers
\(\nu_i\), define
\begin{equation}
  A=\sum_{i=1}^{N}\nu_i,
  \qquad
  D=D_0+D_\eps\eps .
\end{equation}
With the usual projective Feynman-parameter normalization, the parametric
integral has the schematic form
\begin{align}
  I(\nu)
  &=
  \mathcal P(D,\nu)\,
  \frac{\Gamma(A-LD/2)}{\prod_i\Gamma(\nu_i)}
  \int_{x_i\geq0}
  \left[\prod_{i=1}^{N}\dd x_i\,x_i^{\nu_i-1}\right]
  \delta\!\left(1-\sum_i x_i\right)
  \notag\\
  &\hspace{2.8cm}\times
  \U(x)^{A-(L+1)D/2}\,
  \F(x)^{-(A-LD/2)} .
\end{align}
The convention-dependent prefactor \(\mathcal P(D,\nu)\) contains the signs,
powers of \(\ii\pi^{D/2}\), and any global normalization choices.  It is not
part of the sector-local black-box construction.  What the sector processor
needs are the affine powers
\begin{align}
  p_U(\eps)&=A-(L+1)D/2,\\
  p_F(\eps)&=-(A-LD/2),
\end{align}
the parameter weights \(\nu_i-1\), and the retained black-box evaluators for
\(\U\) and \(\F\).

After a sector map \(x=X_s(y)\), the topology-specific sector generator should
separate each possible endpoint monomial source,
\begin{align}
  J_s(y)\prod_i X_{s,i}(y)^{\nu_i-1}N_s(y)
  &=
  J_{s,\mathrm{reg}}(y)
  \prod_a y_a^{j_{s,a}},
  \\
  \U(X_s(y))
  &=
  M_{U,s}(y)\,\psi_s(y),
  \qquad
  M_{U,s}(y)=\prod_a y_a^{u_{s,a}},
  \\
  \F(X_s(y))
  &=
  M_{F,s}(y)\,\phi_s(y),
  \qquad
  M_{F,s}(y)=\prod_a y_a^{f_{s,a}}.
\end{align}
Here \(N_s\) denotes any smooth numerator or extra parameter-space factor.  The
regular functions \(\psi_s\) and \(\phi_s\) are evaluated by black-box
dual-number limits of \(\U\) and \(\F\); neither polynomial is symbolically
opened inside the processor.  The full endpoint power of \(y_a\) is then
\begin{equation}
  \rho_{s,a}(\eps)
  =
  j_{s,a}
  +
  u_{s,a}\,p_U(\eps)
  +
  f_{s,a}\,p_F(\eps),
\end{equation}
with any additional numerator or measure monomial powers included in
\(j_{s,a}\).  This is the object that decides the subtraction problem.  The
implemented examples are logarithmic endpoint cases,
\begin{equation}
  \rho_{s,a}(\eps)=-1+c_{s,a}\eps,
  \qquad
  c_{s,a}\neq0,
\end{equation}
which produce explicit Laurent poles from
\(\int_0^1\dd y_a\,y_a^{-1+c_{s,a}\eps}\).

For a uniform notation, let
\begin{equation}
  P\in\{\U,\F\},
  \qquad
  r_{P,s,a}\in\mathbb N,
  \qquad
  M_{P,s}(y)=\prod_{a\in A_s}y_a^{r_{P,s,a}},
\end{equation}
with \(r_{\U,s,a}=u_{s,a}\), \(r_{\F,s,a}=f_{s,a}\), and \(A_s\) the set of
sector axes on which an endpoint monomial has been declared.  Define
\begin{equation}
  P_s(y)=\widehat P(X_s(y)),
  \qquad
  R_{P,s}(y)=\frac{P_s(y)}{M_{P,s}(y)}.
\end{equation}
Thus \(R_{\U,s}=\psi_s\) and \(R_{\F,s}=\phi_s\).  At an interior point this is
just the quotient above.  At a boundary point, let
\begin{equation}
  B(y)=\{a\in A_s\;|\;y_a=0\},
  \qquad
  \bar B=A_s\setminus B,
\end{equation}
and introduce dual variables \(\eta_b\) for \(b\in B\).  The boundary jet is
formed by
\begin{equation}
  \widehat y_a(\eta)=
  \begin{cases}
    \eta_a, & a\in B,\\
    y_a, & a\notin B .
  \end{cases}
\end{equation}
Then the regular residual is the Taylor coefficient
\begin{equation}
  R_{P,s}(y)
  =
  \frac{
    \left[
      \prod_{b\in B}\eta_b^{r_{P,s,b}}
    \right]\,
    \widehat P\!\left(X_s(\widehat y(\eta))\right)
  }{
    \prod_{a\in \bar B}y_a^{r_{P,s,a}}
  }.
\end{equation}
Here \([\eta^n]\) denotes the Taylor coefficient, not the derivative.  In
derivative notation the same statement is
\begin{equation}
  R_{P,s}(y)
  =
  \frac{
    \left.
    \left(\prod_{b\in B}
      \frac{1}{r_{P,s,b}!}\,
      \partial_{\eta_b}^{r_{P,s,b}}
    \right)
    \widehat P\!\left(X_s(\widehat y(\eta))\right)
    \right|_{\eta=0}
  }{
    \prod_{a\in\bar B}y_a^{r_{P,s,a}}
  }.
\end{equation}
The chain rule is applied at the level of truncated jets:
\begin{equation}
  J^{\mathbf r}_{B}\!\left(P_s\right)(y)
  =
  J^{\mathbf r}_{x}\!\left(\widehat P\right)\!\left(X_s(y)\right)
  \circ
  J^{\mathbf r}_{B}\!\left(X_s\right)(y),
\end{equation}
where
\begin{equation}
  \mathbf r=(r_{P,s,b})_{b\in B}.
\end{equation}
Equivalently, for the first two derivative orders,
\begin{align}
  \partial_a P_s(y)
  &=
  \sum_i
  \frac{\partial \widehat P}{\partial x_i}(X_s(y))\,
  \partial_a X_{s,i}(y),
  \\
  \partial_a\partial_b P_s(y)
  &=
  \sum_{i,j}
  \frac{\partial^2 \widehat P}{\partial x_i\partial x_j}(X_s(y))\,
  \partial_a X_{s,i}(y)\,
  \partial_b X_{s,j}(y)
  \notag\\
  &\quad+
  \sum_i
  \frac{\partial \widehat P}{\partial x_i}(X_s(y))\,
  \partial_a\partial_b X_{s,i}(y).
\end{align}

After extracting the declared monomials, every sector has the form
\begin{equation}
  I_s(\eps)
  =
  \int_{[0,1]^d}\dd^d y\,
  \left[
    \prod_{a\in A_s}y_a^{\rho_{s,a}(\eps)}
  \right]
  G_s(y,\eps),
\end{equation}
where
\begin{equation}
  G_s(y,\eps)
  =
  J_{s,\mathrm{reg}}(y)\,
  R_{\U,s}(y)^{p_U(\eps)}\,
  R_{\F,s}(y)^{p_F(\eps)}
\end{equation}
times any remaining regular numerator factor.  Write
\begin{equation}
  \rho_{s,a}(\eps)=\beta_{s,a}+c_{s,a}\eps,
  \qquad
  \beta_{s,a}\in\mathbb Z,
\end{equation}
and define the endpoint Taylor order
\begin{equation}
  N_{s,a}=\max(-\beta_{s,a}-1,0).
\end{equation}
For one endpoint variable,
\begin{align}
  \int_0^1\dd y\,y^{\beta+c\eps}G(y,\eps)
  &=
  \int_0^1\dd y\,y^{\beta+c\eps}
  \left[
    G(y,\eps)
    -
    \sum_{n=0}^{N}
    \frac{y^n}{n!}
    \partial_y^nG(0,\eps)
  \right]
  \notag\\
  &\quad+
  \sum_{n=0}^{N}
  \frac{\partial_y^nG(0,\eps)}
       {n!\,(\beta+n+1+c\eps)} .
\end{align}
The first line is finite at \(y=0\), while the second line contains the
explicit \(\eps\)-poles.

For several endpoint axes, introduce the Taylor projector
\begin{equation}
  (\mathcal T_a G)(y,\eps)
  =
  \sum_{n=0}^{N_{s,a}}
  \frac{y_a^n}{n!}
  \left[
    \partial_{y_a}^nG(y,\eps)
  \right]_{y_a=0},
  \qquad
  \mathcal R_a=1-\mathcal T_a .
\end{equation}
Then the recursive subtraction identity can be written compactly as
\begin{align}
  I_s(\eps)
  &=
  \sum_{B\subseteq A_s}
  \sum_{\{0\leq n_b\leq N_{s,b}\}_{b\in B}}
  \left[
    \prod_{b\in B}
    \frac{1}{
      n_b!\,
      (\beta_{s,b}+n_b+1+c_{s,b}\eps)
    }
  \right]
  \notag\\
  &\quad\times
  \int
  \dd y_{\bar B}\,
  \left[
    \prod_{a\in A_s\setminus B}
    y_a^{\beta_{s,a}+c_{s,a}\eps}
  \right]
  \left[
    \prod_{a\in A_s\setminus B}\mathcal R_a
  \right]
  \left[
    \left.
    \prod_{b\in B}\partial_{y_b}^{n_b}
    G_s(y,\eps)
    \right|_{y_B=0}
  \right].
\end{align}
In the logarithmic case \(\beta_{s,a}=-1\), one has \(N_{s,a}=0\), and the
identity reduces to
\begin{align}
  I_s(\eps)
  &=
  \sum_{B\subseteq A_s}
  \left[
    \prod_{b\in B}\frac{1}{c_{s,b}\eps}
  \right]
  \int
  \dd y_{\bar B}\,
  \left[
    \prod_{a\in A_s\setminus B}
    y_a^{-1+c_{s,a}\eps}
  \right]
  \notag\\
  &\quad\times
  \left[
    \prod_{a\in A_s\setminus B}(1-\mathcal T_a)
  \right]
  \left[
    G_s(y,\eps)
  \right]_{y_B=0},
  \qquad
  (\mathcal T_aG)(y,\eps)=G(y,\eps)|_{y_a=0}.
\end{align}
Finally, expanding
\begin{equation}
  G_s(y,\eps)=\sum_{k\geq0}\eps^k G_s^{(k)}(y),
  \qquad
  \prod_{a\in A_s\setminus B}y_a^{c_{s,a}\eps}
  =
  \sum_{r\geq0}
  \frac{\eps^r}{r!}
  \left(
    \sum_{a\in A_s\setminus B}c_{s,a}\log y_a
  \right)^r,
\end{equation}
gives the Laurent coefficients explicitly as ordinary finite integrals.

This audit leads to a useful interface boundary.  A topology definition must
retain the global parametric metadata: loop count, propagator powers,
dimension, prefactor description, and affine powers of \(\U\) and \(\F\).  A
sector definition must retain the map, regular sector prefactor, monomial
powers of \(\U\) and \(\F\), monomial powers from the Jacobian/measure or
numerator, and the singular axes.  The generic processor can then assemble
\(\rho_{s,a}\), evaluate \(\psi_s\) and \(\phi_s\) from dualized black-box
evaluators, and apply the appropriate localized subtraction formula.

The current numerical implementation still evaluates only the coefficients
\(\epsilon^{-2},\epsilon^{-1},\epsilon^0\) and implements explicit formulas for
zero, one, and two logarithmic endpoint axes.  Higher-rank endpoint structures
and non-logarithmic endpoint powers require a recursive subtraction engine and
a variable-length Laurent accumulator.  The metadata described above is now
present so such a recursive engine can be added without changing the
black-box/topology/sector boundary.

\section{Outlook: Schwinger sampling with momentum-space numerators}

The construction above assumes that the loop momenta have already been
integrated out, so all non-trivial information enters through the Symanzik
polynomials and possible parameter-space numerator factors.  A complementary
path is useful when the denominator part should still be organized in
Feynman/Schwinger parameters, but the numerator is easier to keep as a
momentum-space black box.  Schematically, start from a regulated integral
\begin{equation}
  I[f]
  =
  \int \dd^{DL}k\,
  \frac{f(k)}
       {\prod_{e\in E}(D_e(k)+\ii0)^{\nu_e}} .
\end{equation}
Using Schwinger parameters, or equivalently the projective Feynman-parameter
form of the same denominator product,
\begin{equation}
  \prod_{e\in E}\frac{1}{(D_e(k)+\ii0)^{\nu_e}}
  =
  \frac{\Gamma(\nu)}{\prod_e\Gamma(\nu_e)}
  \int_{\mathbb P^{|E|-1}_+}
  \prod_e \dd x_e\,x_e^{\nu_e-1}\,
  \left[
    \sum_e x_e(D_e(k)+\ii0)
  \right]^{-\nu},
  \qquad
  \nu=\sum_e\nu_e ,
\end{equation}
one obtains an integral over positive edge parameters and loop momenta,
without replacing \(f(k)\) by tensor-reduction formulae.  The quadratic
denominator combination can be written schematically as
\begin{equation}
  \sum_e x_eD_e(k)
  =
  (k-Q(x))^T A(x)(k-Q(x))+\V(x),
  \qquad
  \V(x)=\frac{\F(x)}{\U(x)} ,
\end{equation}
with the appropriate \(+\ii0\) prescription or later contour deformation
understood.  Introducing an overall Schwinger scale and auxiliary momentum
variables then gives a representation of the schematic form
\begin{align}
  I[f]
  &\propto
  \int_{\mathbb P^{|E|-1}_+}
  \prod_e \dd x_e\,x_e^{\nu_e-1}\,
  \frac{1}{\U(x)^{D/2}\,\V(x)^\omega}
  \notag\\
  &\quad\times
  \int_0^\infty \dd\lambda\,
  \lambda^{\omega-1}e^{-\lambda}
  \int \dd\mu(q)\,
  f\!\left(k(x,q,\lambda)\right),
\end{align}
where
\begin{equation}
  \omega=\sum_e\nu_e-\frac{DL}{2}.
\end{equation}
Here \(\dd\mu(q)\) denotes a fixed reference measure for the auxiliary
momentum variables after the quadratic form has been normalized.  The important
feature is not the particular choice of that reference measure.  The important
feature is that the numerator remains a callable function of momenta, while
the Feynman-parameter part still exposes \(\U\), \(\F\), endpoint monomials,
and sector maps to the same FSD machinery as before.

After sector decomposition of the parameter variables, the intended structure
is
\begin{align}
  I[f]
  &=
  \sum_s
  \int_{[0,1]^{d_s}}\dd y
  \int\dd\lambda\,\dd\mu(q)\,
  \prod_{a\in A_s}y_a^{\rho_{s,a}(\eps)}
  \,
  G_s(y,\lambda,q;\eps)\,
  f\!\left(k_s(y,\lambda,q)\right).
\end{align}
The subtraction operators act on the sector variables \(y_a\).  The auxiliary
variables \(\lambda\) and \(q\) are spectators for the endpoint algebra:
\begin{equation}
  \mathcal S_s
  \left[
    G_s(y,\lambda,q;\eps)\,
    f(k_s(y,\lambda,q))
  \right],
\end{equation}
with \(\mathcal S_s\) the same local Taylor-subtraction operator used for the
ordinary parameter-space integrand.  Thus the future topology object would
store not only \(\U\), \(\F\), and the parametric exponents, but also the
momentum routing, the quadratic denominator data, and the map from sampled
auxiliary variables to loop momenta.  The integrand definition would expose a
momentum-space numerator callback
\begin{equation}
  f:\; k \longmapsto f(k),
\end{equation}
and FSD would provide
\begin{equation}
  z\in[0,1]^n
  \longmapsto
  \left(K_{\mathrm{FSD}}(z),\,k(z)\right),
\end{equation}
so that the numerical integrand is \(K_{\mathrm{FSD}}(z)f(k(z))\).

This is conceptually different from momentum-space tropical sampling, such as
the approach implemented in \texttt{momtrop}; see
\url{https://github.com/alphal00p/momtrop}.  The FSD version would deliberately
use sector decomposition and subtraction for the Feynman-parameter part.  It
therefore inherits the larger sector count and will not compete with tropical
sampling on the factorial-sector problem for finite parameter integrals.  The
advantage is different: the same framework can expose endpoint poles in
\(\eps\), and, once contour deformation or another threshold treatment is
added, it is not tied to a finite Euclidean integration domain.  For generic
higher-order endpoint subtractions one may also need derivatives of the
composed numerator \(f(k_s(y,\lambda,q))\), or an explicit regularity contract
showing that only boundary values are required.  This is a boundary condition
for the future backend, not part of the current scalar implementation.
Moreover, the momentum-space numerator must be understood as a
\(D\)-dimensional numerator, not merely a four-dimensional one.  Rational
terms from the \((D-4)\)-dimensional components have to be retained, for
example through the usual \(\mu\)-term bookkeeping
\begin{equation}
  k_{[D]}^2
  =
  k_{[4]}^2
  +
  \mu^2,
  \qquad
  \mu^2 \equiv k_{[D-4]}^2 ,
\end{equation}
with signs and metric conventions fixed consistently with the chosen
denominator prescription.


\section{Implementation in \texttt{FSD}}

The repository contains the modular implementation used for the numerical
checks described above.  It is referred to simply as \texttt{FSD}.  The
executable entry point is
\begin{equation}
  \texttt{FSD.py}.
\end{equation}
A local virtual environment is created at
\begin{equation}
  \texttt{.venv},
\end{equation}
and OneLOopBridge is required as an external dependency.  The benchmark values
are obtained from the Python binding calls
\begin{align}
  &\texttt{oneloop\_bridge.three\_point(0,0,s,m2,m2,m2)},
  \\
  &\texttt{oneloop\_bridge.four\_point(0,0,0,0,s12,s23,m2,m2,m2,m2)}.
\end{align}
The only implementation described here is the modular \texttt{FSD} code.

\subsection{Code organization}

The implementation is organized so that the black-box rule is visible directly
in the code layout:
\begin{itemize}
  \item \texttt{FSD.py} is the command-line steering layer.  It parses options,
  validates kinematics, builds the request object, prints summaries, calls the
  benchmark, launches the integration, and renders the final table.
  \item \texttt{sectors\_generator.py} owns the declarative sector data:
  sector maps \(X_s\), regular sector prefactors, endpoint monomials
  \(M_{U,s}\) and \(M_{F,s}\), monomial powers from Jacobian/measure or
  numerator factors, singular axes, and generated Symbolica callbacks for
  these expressions.
  \item \texttt{integrand.py} owns the topology definition and the generic
  \texttt{SectorProcessor}.  The topology retains \(\U\) and \(\F\) for
  display and evaluator construction, but the processor never substitutes a
  sector map into them symbolically.
  \item \texttt{integrator.py} owns the Havana sampling loop, manual Laurent
  coefficient accumulation, progress reporting, target-accuracy stopping, and
  hot-path timing profile.
  \item \texttt{benchmark.py}, \texttt{formatting.py}, \texttt{definitions.py},
  and \texttt{utils.py} isolate OneLOopBridge access, colored output, shared
  data containers, and parenthesis uncertainty formatting.
  \item \texttt{dot\_topology.py} is the scaffold for future
  GammaLoop-convention DOT-file input.  It is wired into the same topology and
  sector interfaces, but its parser, topology builder, sector issuer, and
  benchmark mapping intentionally raise \texttt{NotImplementedError}.
\end{itemize}

\subsection{DOT-file extension path}

The command-line option
\begin{equation}
  \texttt{--dot-file path/to/integral.dot}
\end{equation}
selects the future GammaLoop DOT-backed topology path.  The current code
validates that the file exists and then dispatches to
\texttt{GammaLoopDotTopologyBuilder}.  The intended development sequence is
\begin{align}
  \texttt{DOT file}
  &\longrightarrow
  \texttt{GammaLoopTopologyData},
  \\
  \texttt{GammaLoopTopologyData}
  &\longrightarrow
  \texttt{TopologyDefinition},
  \\
  \texttt{GammaLoopTopologyData}
  &\longrightarrow
  \left\{\texttt{SectorDefinition}\right\}.
\end{align}
This keeps arbitrary future topologies on the same path as the built-in
triangle and box examples: topology-specific information is converted into
retained Symanzik expressions and declarative sector metadata, while the
\texttt{SectorProcessor} remains generic and continues to consume only
black-box evaluators and sector callbacks.

\subsection{Prepared sectors and black-box processing}

The code mirrors the split used in this note.  Sectors are prepared
independently of the Symanzik polynomials.  Each sector stores only
\begin{equation}
  \left\{
    X_s,\,
    J_{s,\mathrm{reg}},\,
    M_{U,s},\,
    M_{F,s},\,
    \text{measure/numerator monomial powers},\,
    \text{singular axes}
  \right\},
\end{equation}
plus Symbolica evaluators generated from these known sector expressions.  The
processing layer receives these prepared sectors together with black-box
Symbolica evaluators
\begin{equation}
  \hU(x;p),
  \qquad
  \hF(x;p),
\end{equation}
where \(p\) denotes the external invariants and masses.  After the evaluators
are built, the processing code does not inspect or symbolically rewrite the
expressions for \(\U\) or \(\F\).

At runtime the processor receives only a \texttt{SectorDefinition}, a
\texttt{TopologyDefinition}, and numeric sample points.  For a sector point
\(y\) it evaluates
\begin{equation}
  x = X_s(y),
  \qquad
  \hU(x),
  \qquad
  \hF(x),
  \qquad
  J_s(y),
\end{equation}
through prebuilt Symbolica evaluators.  If the sector has no endpoint
monomial, the contribution is evaluated directly.  If a monomial is declared,
the residuals
\begin{equation}
  \psi_s(y)=\frac{\hU(X_s(y))}{M_{U,s}(y)},
  \qquad
  \phi_s(y)=\frac{\hF(X_s(y))}{M_{F,s}(y)}
\end{equation}
are evaluated by these quotients in the interior.  On endpoint samples the
same quantities are recovered from Taylor coefficients of dualized black-box
\(\hU\) and \(\hF\) evaluators.  The dual jets of the known sector map provide
the chain-rule input, so no symbolic expansion of \(\U(X_s(y))\) or
\(\F(X_s(y))\) is performed.  The scalar \(\hU\) and \(\hF\) evaluators are
constructed once with the topology.  Dualized evaluators are cached by their
requested Taylor shape and are obtained by copying these scalar evaluators
before applying Symbolica's dualization step.

The generic processor then constructs
\begin{equation}
  g_s(y;\eps)
  =
  J_{s,\mathrm{reg}}(y)\,
  \psi_s(y)^{p_U(\eps)}\,
  \phi_s(y)^{p_F(\eps)},
\end{equation}
assembles the endpoint powers \(\rho_{s,a}(\eps)\), and applies the declared
localized subtraction formula.  Zero, one, and two logarithmic singular axes
are supported explicitly.  Higher-rank or non-logarithmic endpoint structures
raise an error instead of being guessed.

\subsection{Havana sampling and stopping}

\texttt{FSD} uses Symbolica's Havana integrator as a sampler rather than as a
single black-box scalar integrator.  One discrete dimension chooses the sector
and each sector has a continuous adaptive grid of the same dimension:
\begin{equation}
  \text{discrete sector id}
  \quad\otimes\quad
  y\in[0,1]^d.
\end{equation}
Lower-support subtraction pieces are localized into the same highest-dimensional
sample space by dummy unit-density dependence on unused variables.  This avoids
separate grids for bulk, edge, and corner pieces.

Each sampled point contributes to all Laurent coefficients through the same
manual accumulator,
\begin{equation}
  A_k
  =
  \frac{1}{N}\sum_{n=1}^{N}
  w_n\,I_{k,s_n}(y_n),
  \qquad
  k\in\{-2,-1,0\}.
\end{equation}
Havana is trained only on the finite-part magnitude used as an importance
function.  The adaptive grid is cloned at the beginning of an iteration; batch
training samples are accumulated into the clone and merged back into the live
grid only at the end of a completed iteration.

The target-accuracy stop is evaluated after each accumulated batch, not only
after a full Havana iteration.  Thus the final returned coefficients contain
all completed batches up to the stopping point.  In parallel mode the driver
keeps only a worker-sized window of batches in flight, so early stopping does
not submit the entire iteration unnecessarily.  The option
\(\texttt{--batch-size}\) therefore controls both vectorization size and
mid-iteration stopping granularity.

\subsection{Timing profile}

The timing profile reported by the progress bar and final table is a work-time
profile:
\begin{equation}
  T_{\rm prof}
  =
  T_{\rm Python}
  +
  T_{\rm eval}
  +
  T_{\rm Havana}.
\end{equation}
Here \(T_{\rm eval}\) is the sum of Symbolica evaluator wall times over all
workers, \(T_{\rm Python}\) is NumPy and Python glue on workers plus
master-side submission, result retrieval, accumulation, and progress rendering,
and \(T_{\rm Havana}\) is Havana sampling, training accumulation, grid merge,
and grid update time.  The displayed evaluator average is
\begin{equation}
  \langle t_{\rm eval}\rangle
  =
  \frac{10^6\,T_{\rm eval}}{N_{\rm accepted}}
  \quad
  \text{microseconds per sample per worker},
\end{equation}
because \(T_{\rm eval}\) is already summed over workers.

\subsection{Implemented modes and conventions}

For the triangle, the command-line interface supports massive and massless
modes.  In the validated finite massive region, endpoint monomial
factorization is not applied.  The implementation evaluates
\begin{equation}
  C_0(s;m^2)
  =
  -
  \sum_{a=1}^{2}
  \int_0^1\dd t
  \int_0^1\dd z\,
  J_{S_a}(t,z)\,
  \hU(X_{S_a}(t,z))^{-1}\,
  \hF(X_{S_a}(t,z))^{-1}.
\end{equation}
In massless Euclidean mode, the implementation uses the subtraction formula of
Section~3.2 and returns the coefficients of
\begin{equation}
  C_0(s;0)
  =
  \frac{C_{-2}}{\epsilon^2}
  +
  \frac{C_{-1}}{\epsilon}
  +
  C_0
  +
  \mathcal O(\epsilon).
\end{equation}

The default output convention is the OneLOopBridge convention.  In this
convention the universal Gamma/Euler factors are stripped from the displayed
Laurent coefficients.  For massless endpoint representations this also amounts
to applying the finite convention shift described in the direct construction.
The option \(\texttt{--prefactor-convention}\) controls only the overall
scalar-integral normalization shown in the output table.  The default value
\(\texttt{raw}\) displays the raw OneLOopBridge normalization.  The value
\(\texttt{feynman}\) multiplies both the \texttt{FSD} coefficients and the
benchmark coefficients by
\begin{equation}
  \texttt{TO\_FEYNMAN}
  =
  -\frac{1}{16\pi^2}.
\end{equation}
In either display convention, values with Monte Carlo uncertainty are printed
with the two-significant-digit parenthesis notation.

\subsection{Implemented box example}

The implementation also contains a scalar box example with massless external
legs,
\begin{equation}
  D_0(0,0,0,0,s_{12},s_{23};m^2).
\end{equation}
For equal internal masses, the simplex polynomial used by the implementation is
\begin{equation}
  \U=x_0+x_1+x_2+x_3,
  \qquad
  \F
  =
  m^2\U^2
  -
  s_{12}x_0x_2
  -
  s_{23}x_1x_3
  -
  \ii0.
\end{equation}
After imposing \(\U=1\), the finite four-point integral in \(D=4\) is
\begin{equation}
  D_0
  =
  \sum_{b=0}^{3}
  \int_0^1\dd y_0
  \int_0^1\dd y_1
  \int_0^1\dd y_2\,
  J_b(y)\,
  \hF(X_b(y))^{-2}.
\end{equation}
Here \(b\) labels the primary sector in which \(x_b\) is the largest Feynman
parameter.  In sector \(b\),
\begin{equation}
  x_b=\frac{1}{1+y_0+y_1+y_2},
  \qquad
  x_{j\neq b}=\frac{y_j}{1+y_0+y_1+y_2},
\end{equation}
with the three \(y_j\)'s assigned in increasing order to the non-dominant
parameters.  The Jacobian is
\begin{equation}
  J_b(y)=\frac{1}{(1+y_0+y_1+y_2)^4}.
\end{equation}

For \(m=0\), each primary sector has the form
\begin{equation}
  \F(X_b(y))
  =
  \frac{A_b\,y_\ell+B_b\,y_r y_s}
       {(1+y_0+y_1+y_2)^2},
\end{equation}
where one of the two bilinears \(x_0x_2\) and \(x_1x_3\) contains the dominant
parameter and is therefore linear in the primary variables.  The implementation
does not need the symbolic values of \(A_b\) and \(B_b\); it only uses the known
positions of the linear variable \(y_\ell\) and product variables \(y_r,y_s\).
The three secondary sectors are
\begin{align}
  L:\quad& y_r=uv,\quad y_\ell=u,\quad y_s=w,
  &
  M_L&=u,
  \\
  P:\quad& y_r=u,\quad y_s=v,\quad y_\ell=uvw,
  &
  M_P&=uv,
  \\
  Q:\quad& y_r=u,\quad y_\ell=uv,\quad y_s=vw,
  &
  M_Q&=uv.
\end{align}
The secondary Jacobian equals the displayed monomial in each case, so the
regular prefactor after monomial extraction is again
\begin{equation}
  \frac{1}{(1+y_0(u,v,w)+y_1(u,v,w)+y_2(u,v,w))^4}.
\end{equation}
The residual coefficient
\begin{equation}
  \phi_{b\alpha}(u,v,w)
  =
  \frac{\F(X_{b\alpha}(u,v,w))}{M_\alpha(u,v,w)}
\end{equation}
is evaluated by the black-box \(\F\) away from the endpoint and by dual-number
Taylor coefficients at \(u=0\) and \(v=0\).  The one-monomial sector \(L\) uses
the one-variable subtraction formula, while the \(P\) and \(Q\) sectors use the
two-variable subtraction formula.

The raw coefficients produced directly by the massless box sectors are
converted to the OneLOopBridge stripped convention as
\begin{equation}
  D_{-2}=A_{-2},
  \qquad
  D_{-1}=A_{-1}+A_{-2},
  \qquad
  D_0=A_0+A_{-1}+\frac{\pi^2}{6}A_{-2}.
\end{equation}
The current massless box command validates Euclidean kinematics,
\(s_{12}<0\) and \(s_{23}<0\).  Non-Euclidean massless kinematics are outside
the present implementation.

\subsection{Example commands}

The massive triangle check can be run with
\begin{equation}
  \texttt{.venv/bin/python FSD.py --integral triangle --s 1.0 --m 1.0}.
\end{equation}
The same check in the Feynman-normalized prefactor convention can be displayed
with
\begin{align}
  &\texttt{.venv/bin/python FSD.py --integral triangle --s 1.0 --m 1.0}
  \notag\\
  &\texttt{--prefactor-convention feynman}.
\end{align}
The massless Euclidean triangle check can be run with
\begin{equation}
  \texttt{.venv/bin/python FSD.py --integral triangle --s -1.0 --m 0.0}.
\end{equation}
The massive box example can be run with
\begin{equation}
  \texttt{.venv/bin/python FSD.py --integral box --s12 0.5 --s23 0.7 --m 1.0}.
\end{equation}
The massless Euclidean box check can be run with
\begin{equation}
  \texttt{.venv/bin/python FSD.py --integral box --s12 -1.0 --s23 -2.0 --m 0.0}.
\end{equation}
Massless timelike commands, such as
\(\texttt{--integral triangle --s 1.0 --m 0.0}\), are rejected by validation.
OneLOopBridge is required in all runs because the implementation always
compares the Laurent coefficients to a benchmark value.  The implementation
intentionally imports neither \texttt{scipy} nor \texttt{sympy}; all symbolic
handling is done through Symbolica evaluators.

\end{document}
